\documentclass[10pt]{article}
\usepackage[T1]{fontenc}
\usepackage[utf8]{inputenc}
\usepackage[french]{babel}
\usepackage{newpxtext,newpxmath}
\usepackage[a4paper,top=1.5cm,bottom=1.35cm,left=1.65cm,right=1.65cm]{geometry}
\usepackage{microtype}
\usepackage{xcolor}
\usepackage{titlesec}
\usepackage[most]{tcolorbox}
\usepackage{enumitem}
\usepackage{array}
\usepackage{booktabs}
\usepackage{fancyhdr}

\definecolor{navy}{HTML}{1F3A5F}
\definecolor{grey}{HTML}{54606E}
\definecolor{green}{HTML}{2C6E49}
\definecolor{amber}{HTML}{9C6A16}
\definecolor{red}{HTML}{A6322E}
\definecolor{plum}{HTML}{7A4EA8}
\definecolor{boxgray}{HTML}{F3F5F8}
\definecolor{boxred}{HTML}{FCF2F1}
\definecolor{boxgreen}{HTML}{F1F7F3}
\definecolor{boxplum}{HTML}{F7F3FB}
\definecolor{boxedge}{HTML}{D9E0EA}

\setlength{\parindent}{0pt}
\setlength{\parskip}{2.6pt}
\linespread{1.0}

\pagestyle{fancy}\fancyhf{}
\lhead{\small\color{grey}Scella}
\rhead{\small\color{grey}Couples -- classement intrinsèque -- 11 août 2026}
\cfoot{\small\color{grey}\thepage}
\renewcommand{\headrulewidth}{0.3pt}

\titleformat{\section}{\color{navy}\large\bfseries}{}{0pt}{}
\titlespacing{\section}{0pt}{9pt}{3pt}
\titleformat{\subsection}{\color{navy}\bfseries}{}{0pt}{}
\titlespacing{\subsection}{0pt}{7pt}{1.5pt}

\newtcolorbox{these}{colback=boxgray,colframe=boxedge,boxrule=0.5pt,arc=2.5pt,
  left=8pt,right=8pt,top=4.5pt,bottom=4.5pt}
\newtcolorbox{alerte}{colback=boxred,colframe=red!55,boxrule=0.8pt,arc=2.5pt,
  left=8pt,right=8pt,top=4.5pt,bottom=4.5pt}
\newtcolorbox{acquis}{colback=boxgreen,colframe=green!45,boxrule=0.8pt,arc=2.5pt,
  left=8pt,right=8pt,top=4.5pt,bottom=4.5pt}
\newtcolorbox{piste}{colback=boxplum,colframe=plum!45,boxrule=0.8pt,arc=2.5pt,
  left=8pt,right=8pt,top=4.5pt,bottom=4.5pt}

\newcommand{\best}[1]{{\color{green}\textbf{#1}}}
\newcommand{\pire}[1]{{\color{red}\textbf{#1}}}
\newcommand{\sz}[1]{{\small\color{grey}\textsf{#1}}}
\newcommand{\code}[1]{\texttt{\small #1}}
\newcommand{\x}{$\times$}
\newcommand{\hausse}{{\color{green}$\uparrow$}}
\newcommand{\baisse}{{\color{red}$\downarrow$}}

\begin{document}

\begin{center}
{\Huge\bfseries\color{navy} Scella}\\[4pt]
{\large\color{grey}\itshape Les couples, reclassés sans finalité extérieure}\\[4pt]
{\small\color{grey} 11 août 2026 \; · \; après l'abandon de la piste Kovant \; · \; quatre lentilles intrinsèques}
\end{center}
\vspace{2pt}

\begin{these}
\textbf{Le résultat central, et il contredit le précédent.} Le 5 août, retirer la lentille d'objectif
au niveau des \emph{composants} ne changeait presque rien -- un seul échange entre deux ex \ae quo. Au
niveau des \textbf{couples}, la retirer change beaucoup. Les \textbf{vérifieurs zéro-dépendance
passent premiers} alors qu'ils n'étaient \emph{même pas classés} ; \textbf{PII \x{} scellement perd
trois rangs} ; et la \textbf{co-signature R3} -- que j'avais désignée la veille comme le couple le plus
important du dépôt -- \textbf{tombe septième}. Aucun des deux exercices n'était mal fait : la lentille
extérieure filtrait simplement \emph{plus étroitement} sur des propriétés composées que sur des
composants.
\end{these}

\section{1 · Ce qui a été retiré, et ce que cela laisse}

Le classement des couples du 10 août notait sur quatre axes, dont le quatrième était \emph{« valeur
dans la conversation Kovant »}. Cet axe est mort le 11 août avec la piste elle-même. Plutôt que d'en
inventer un -- et surtout plutôt que de lui substituer en silence une autre finalité non déclarée --
ce document reprend le cadre \textbf{déjà établi le 5 août} dans le \emph{Classement du potentiel
intrinsèque} : la \textbf{maturité} rapportée à l'effort restant, la \textbf{rareté} pure,
l'\textbf{adversaire} en revue diligente, et la \textbf{démontrabilité} devant un ingénieur sceptique.

Ces quatre lentilles ne présupposent aucun objectif. Elles répondent à une question plus simple et
plus dure : \emph{qu'est-ce qui vaut le plus, indépendamment de ce qu'on veut en faire ?}

La substitution n'est pas neutre, et c'est là que se joue tout le reclassement. Trois de mes anciens
axes se projettent proprement -- la \emph{distance au servi} devient la maturité, le \emph{mordant}
devient l'adversaire, et la démontrabilité existait déjà. Mais l'axe Kovant est remplacé par la
\textbf{rareté}, qui ne mesure pas du tout la même chose. « Un acheteur conformité pose-t-il cette
question ? » et « cette paire existe-t-elle ailleurs ? » sont deux tests orthogonaux -- et plusieurs
couples qui réussissaient le premier échouent au second.

\begin{alerte}
\textbf{Ce que ce document mesure, et ce qu'il juge.} Les rapports précédents affirmaient des faits
vérifiables : « 515 tests, 0 échec », « \code{remote\_participant.rs:404} renvoie \code{Stub} ». Ce
document-ci fait autre chose. Les \emph{couples} et leurs \emph{emplacements} restent des faits, et ils
sont inchangés. Mais les \textbf{notes} ci-dessous sont un \textbf{jugement}, structuré par des
critères explicites pour être discutable -- pas une mesure. Un désaccord sur « rareté = 3 » est
légitime et ne se tranche par aucune commande. La distinction est faite ici plutôt que laissée au
lecteur, parce que confondre les deux régimes serait exactement le défaut que ce dépôt combat.
\end{alerte}

\section{2 · Le classement intrinsèque}

\renewcommand{\arraystretch}{1.2}
\begin{center}\small
\begin{tabular}{@{}r >{\raggedright\arraybackslash}p{5.1cm} cccc c c c >{\raggedright\arraybackslash}p{4.4cm}@{}}
\toprule
& \textbf{Le couple} & \sz{Mat} & \sz{Rar} & \sz{Adv} & \sz{Dém} & \textbf{Moy} & \sz{Éc.} & \sz{$\Delta$} & \textbf{Ce qui le porte / le plombe} \\
\midrule
\textbf{1} & \textbf{Les vérifieurs zéro-dépendance} & \best{1} & 3 & \best{1} & 3 & \textbf{2,00} & 2 & \hausse\sz{h.c.} &
  Le seul sans aucun caveat ; tenu \emph{deux} fois \\
\textbf{2} & \textbf{Topologie déclarée \x{} reçu injectif} & 2 & 4 & 2 & 2 & \textbf{2,50} & \best{2} & \baisse\sz{1} &
  Couvre le cas qu'une implémentation naïve rate \\
\textbf{3} & \textbf{Reçu v4 \x{} ancre 3a} & 3 & 5 & 3 & \best{1} & \textbf{3,00} & 4 & \sz{=} &
  Se démontre par un vert gênant ; ancre détenue \\
\textbf{4} & \textbf{Canari H1 \x{} participants} & 6 & 2 & 4 & 4 & \textbf{4,00} & 4 & \hausse\sz{1} &
  Réfute sa propre note ; échappatoire déclarée \\
5\,= & Caviardage PII \x{} scellement & 4 & 7 & 6 & 6 & 5,75 & 3 & \pire{\baisse\,3} &
  L'ordre est juste ; sans policy, sortie en clair \\
5\,= & Reçu \x{} coût réel par appel & 5 & \pire{8} & 5 & 5 & 5,75 & 3 & \baisse\sz{1} &
  Solide et borné ; catégorie très occupée \\
7 & \textbf{Co-signature R3 \x{} HMAC} & \pire{8} & \best{1} & \pire{8} & \pire{8} & 6,25 & \pire{7} & \sz{=} &
  Conception supérieure ; \textbf{ne se déclenche jamais} \\
8 & \code{slm\_mediated} \x{} sa chute & 7 & 6 & 7 & 7 & 6,75 & \best{1} & \baisse\sz{2} &
  Patron élégant ; portée étroite, voisin d'une fissure \\
\bottomrule
\end{tabular}
\end{center}

\sz{$\Delta$ = déplacement par rapport au classement du 10 août, qui incluait la lentille Kovant.
« h.c. » = hors concours : les vérifieurs y étaient rangés en « amplificateur », donc non classés.
Éc. = amplitude entre le meilleur et le pire rang obtenus par un même couple.}

\section{3 · Les trois déplacements qui comptent}

\subsection{3.1 · Les vérifieurs zéro-dépendance passent premiers -- et ils n'étaient pas classés}

C'est le déplacement le plus instructif, parce qu'il porte sur une décision que j'avais prise
\emph{avant} de noter. Je les avais rangés hors concours, en « amplificateur », au motif qu'ils ne
produisent aucune propriété neuve : ils rendent les autres croyables. C'était un jugement
\textbf{avec acheteur}. Un vérifieur ne se vend pas -- il ne fait que permettre de contrôler ce qu'on
vend.

Sans cible extérieure, ce raisonnement s'inverse : \textbf{permettre de vérifier n'est plus un
préalable à la valeur, c'est la valeur}. Et une fois admis en concours, ils gagnent sans discussion.
Ils sont \best{1} en maturité -- les seuls du lot auxquels il ne manque rien et dont aucun caveat n'a
été relevé. Ils sont \best{1} devant l'adversaire, parce que la revendication se contrôle en une ligne
(\code{cargo tree -e normal}) et qu'aucune rhétorique n'y résiste. Et leur rareté est haute : atteindre
zéro dépendance impose un SHA-256 et un HMAC écrits à la main, discipline que presque personne ne tient
-- et qui a été tenue \textbf{deux fois}, sur deux piles indépendantes. Ce n'est pas un accident de
graphe de dépendances, c'est une intention répétée.

\subsection{3.2 · PII \x{} scellement perd trois rangs -- la plus grosse perte, et elle est méritée}

Il était deuxième. Son rang était porté par le RGPD et l'achat conformité, c'est-à-dire par
\emph{la question d'un acheteur} -- pas par une propriété. C'était l'axe le plus dépendant de la cible,
et il disparaît avec elle.

Ce qui reste est réel mais plus modeste. L'ordre est juste, et c'est ce qui en fait un vrai couple :
le caviardage précède tout usage du texte, donc le sceau couvre ce qui est réellement sorti. Mais la
\textbf{rareté est faible} (\sz{7/8}) -- caviarder des données personnelles avant un appel tiers est
une pratique répandue, et seule la \emph{discipline d'ordre} est distinctive. Et l'adversaire y trouve
prise (\sz{6/8}) : sans policy tenant, le prompt sort en clair. Le code le documente au lieu de le
taire, ce qui est à son honneur, mais un examinateur hostile s'y arrête et il a raison de s'y arrêter.

\subsection{3.3 · La co-signature R3 tombe septième -- le résultat inconfortable}

La veille, je vous l'avais présentée comme le couple le plus important du dépôt. Cette phrase était
adossée à une question qui n'existe plus : \emph{qu'est-ce que cela nous autorise à dire ?} La
non-répudiation avait un prix élevé parce que les \emph{affirmations} étaient l'enjeu.

Sans cet enjeu, R3 s'effondre. Il est \best{1} en rareté -- la conception est réellement supérieure :
HMAC tenant et Ed25519 du n\oe ud sont orthogonaux, donc \emph{forger un R3 exige les deux clés}, et la
règle de non-élévation redescend automatiquement toute co-signature invérifiable sans jamais se taire.
Mais il est \pire{8} sur les trois autres lentilles, et pour une seule raison : \textbf{aucun
producteur}. Le registre de n\oe uds est vide en production, rien ne co-signe, et le barreau le plus
faible est R4 par construction. Devant un adversaire, cela se tue en une phrase. Devant un ingénieur
sceptique, il n'y a rien à montrer. \textbf{Une valeur réalisée nulle ne se rattrape pas par de la
rareté.}

\begin{piste}
\textbf{Mais son écart de \pire{7} est de loin le plus grand du tableau}, et c'est là qu'est
l'information. Un couple qui obtient le premier rang sur une lentille et le dernier sur trois autres
n'est pas « moyen » : il est \textbf{instable}. La formule juste est \emph{pire réalisé, meilleur
latent}. Un seul fil -- un producteur de co-signature sur le chemin servi -- le ferait passer de
septième à disputer la première place. Ma phrase de la veille était juste sur le potentiel et fausse
sur l'état ; le classement les sépare enfin.
\end{piste}

\section{4 · Ce que ce reclassement dit de la méthode}

Le document du 5 août concluait, avec raison, qu'enlever la lentille d'embauche \textbf{ne changeait
presque rien} -- un seul échange entre deux ex \ae quo -- et en tirait que le classement mesurait bien
la valeur intrinsèque. Ici, l'inverse : un couple entre en tête depuis le hors-concours, un autre perd
trois rangs, et celui que j'avais désigné comme le plus important tombe septième.

Les deux résultats sont compatibles, et leur écart est lui-même le résultat le plus utile de ce
document.

Un \textbf{composant} vaut à peu près ce qu'il vaut : sa maturité, sa rareté et sa solidité ne
dépendent guère de qui l'examine. Une \textbf{propriété composée}, elle, ne se juge pas dans l'absolu
-- elle se juge sur \emph{qui pose la question à laquelle elle répond}. Changer le questionneur change
donc la réponse, et d'autant plus que la question d'origine était spécifique. « Un acheteur conformité
en achats pose-t-il cette question ? » est un filtre bien plus étroit que « cette paire est-elle rare
et solide ? ».

\begin{acquis}
\textbf{La règle réutilisable.} Quand on retire l'axe extérieur d'un classement, ne pas supposer la
stabilité : \emph{vérifier} si cet axe filtrait plus étroitement que les axes intrinsèques. S'il
filtrait plus large, rien ne bouge et le classement était déjà intrinsèque. S'il filtrait plus étroit
-- le cas ici -- il portait des rangs qui ne tenaient qu'à lui, et le reclassement est obligatoire.
\end{acquis}

\section{5 · Conséquences pratiques}

\textbf{Le rapport du 10 août a une hiérarchie périmée, pas des faits périmés.}
\code{Scella\_Meilleurs\_Couples\_2026-08-11.tex} note sur la lentille Kovant et recommande « montrer
le rang 1, poser le fil du rang 7 ». Tout ce qu'il \emph{établit} reste vrai : les couples, les faux
couples, les emplacements, l'annexe rejouable, et la correction en cours d'écriture qui a fait passer
R3 de « à inventer » à « à un fil ». C'est son \textbf{ordre} qui ne vaut plus. Les deux documents se
lisent donc ensemble : celui-là pour les faits, celui-ci pour la hiérarchie.

\textbf{Les faux couples ne bougent pas, et c'est normal.} Fil B \x{} souveraineté, \code{config\_hash}
\x{} ce qui a tourné, \mbox{« vrais SLM »} \x{} essaim : ce sont des constats de \emph{fausseté}, pas
des notes. Ils ne dépendaient d'aucune finalité et n'en dépendent toujours pas. Leur seule évolution est
qu'ils ne risquent plus d'être prononcés devant Kovant -- ils risquent d'être prononcés devant
quelqu'un d'autre.

\textbf{Aucune cible de remplacement n'est déclarée, et ce document n'en invente pas.} Le classement
ci-dessus est sans finalité, ce qui est un état parfaitement tenable -- c'était déjà la posture du
5 août. Si une nouvelle cible apparaît, la quatrième lentille se rejoue et le classement doit être
refait, en sachant maintenant qu'il bougera.

\section{6 · La sensibilité du classement -- ce qui le ferait bouger}

Un classement fondé sur un jugement doit dire à quoi il est sensible, sinon il se lit comme un verdict
alors qu'il n'est qu'un instantané. Voici, pour chaque couple, l'événement unique qui déplacerait son
rang -- et le sens du déplacement. C'est aussi, accessoirement, la liste des chantiers rangés par
effet plutôt que par envie.

\renewcommand{\arraystretch}{1.2}
\begin{center}\small
\begin{tabular}{@{}r >{\raggedright\arraybackslash}p{4.6cm} >{\raggedright\arraybackslash}p{7.1cm} >{\raggedright\arraybackslash}p{3.5cm}@{}}
\toprule
& \textbf{Le couple} & \textbf{L'événement unique qui le déplacerait} & \textbf{Effet} \\
\midrule
1 & Vérifieurs zéro-dép & \emph{aucun} qui le monte -- il est au plafond. En revanche une seule
  dépendance ajoutée le ferait chuter de partout & \sz{ne peut que} \baisse \\
2 & Topologie \x{} reçu injectif & \textbf{M4} : un \code{config\_hash} \emph{rapporté} et non plus
  déclaré fermerait le faux couple F2 qui le borde & \hausse{} \sz{vers la 1\textsuperscript{re}} \\
3 & Reçu v4 \x{} ancre & \textbf{3b} : une ancre tierce, l'ôtant des mains de l'opérateur --
  \code{verify\_anchor} existe déjà & \hausse{} \sz{sur l'adversaire} \\
4 & Canari H1 \x{} participants & retirer l'échappatoire \code{h1\_canary.mode = "off"}, ou la rendre
  refusable par configuration & \hausse{} \sz{modeste} \\
5\,= & PII \x{} scellement & une policy tenant par \textbf{défaut fail-closed} : la sortie en clair
  cesserait d'être le comportement d'un tenant non configuré & \hausse{} \sz{2 rangs} \\
5\,= & Reçu \x{} coût & rien de bon marché -- sa faiblesse est la \emph{rareté}, qu'aucun travail
  interne ne change & \sz{stable} \\
\textbf{7} & \textbf{Co-signature R3} & \textbf{un producteur} sur le chemin servi, plus un registre
  de n\oe uds peuplé en production & \best{\hausse{} vers la 1\textsuperscript{re}--2\textsuperscript{e}} \\
8 & \code{slm\_mediated} & fermer la fissure F1 : contraindre l'URL du SLM et décompter son egress
  & \hausse{} \sz{modeste} \\
\bottomrule
\end{tabular}
\end{center}

Deux lectures de ce tableau, et la seconde est la plus utile.

La première est arithmétique : \textbf{un seul chantier déplace un couple de six rangs}, et c'est le
producteur de co-signature. Tous les autres valent un ou deux rangs. Si l'on classe les travaux par
\emph{effet sur la valeur} plutôt que par difficulté ressentie, la hiérarchie est donc écrasée par un
seul item -- ce qui est exactement ce que mesurait l'écart de \pire{7} du tableau précédent.

La seconde est plus dérangeante : \textbf{le couple classé premier ne peut que descendre}. Les
vérifieurs zéro-dépendance n'ont aucune marge de progression -- ils sont au plafond sur trois lentilles
sur quatre -- et une seule dépendance ajoutée, un jour, par commodité, les ferait chuter de partout à
la fois.

\begin{alerte}
\textbf{Et rien ne les protège. C'est un fait, pas un jugement, et il a été vérifié en écrivant cette
section.} La contrainte zéro-dépendance est \emph{documentée} en commentaire dans quatre
\code{Cargo.toml}, et sa \emph{fidélité} est prouvée par des vecteurs golden et un cross-check. Mais
l'invariant lui-même -- « \code{cargo tree -e normal} ne montre que ce crate » -- n'est gardé par
\textbf{aucun test et aucun job d'intégration continue}. Le seul contrôle exécutable du dépôt vit dans
\code{scripts/demo-consensus-receipt.sh}, un script de démonstration que rien ne déclenche
automatiquement. L'actif le mieux classé du portefeuille est donc protégé par une ligne que personne
n'exécute, et sa dégradation serait \textbf{silencieuse} : un \code{cargo add} passerait la revue sans
qu'aucun voyant ne s'allume. \sz{Correctif : porter ce contrôle en test, taille XS.}
\end{alerte}

\begin{these}
\small\textbf{Portée et statut.} Les couples, leurs emplacements et les faux couples sont des faits
établis contre le code du 11 août et re-vérifiables par l'annexe du rapport compagnon. Les
\textbf{notes} de ce document-ci sont un jugement structuré, pas une mesure. La couche nano et
\code{efe-gold-test-gate} n'ont été examinés dans aucun des deux. Et le rappel qui ne dépend d'aucune
cible : le reçu atteste l'\textbf{intégrité} d'une suite de pas, jamais la \textbf{correction} d'une
décision.
\end{these}

\end{document}
